% Options for packages loaded elsewhere
\PassOptionsToPackage{unicode}{hyperref}
\PassOptionsToPackage{hyphens}{url}
\PassOptionsToPackage{dvipsnames,svgnames,x11names}{xcolor}
%
\documentclass[
  11pt,
]{article}
\usepackage{amsmath,amssymb}
\usepackage{lmodern}
\usepackage{iftex}
\ifPDFTeX
  \usepackage[T1]{fontenc}
  \usepackage[utf8]{inputenc}
  \usepackage{textcomp} % provide euro and other symbols
\else % if luatex or xetex
  \usepackage{unicode-math}
  \defaultfontfeatures{Scale=MatchLowercase}
  \defaultfontfeatures[\rmfamily]{Ligatures=TeX,Scale=1}
\fi
% Use upquote if available, for straight quotes in verbatim environments
\IfFileExists{upquote.sty}{\usepackage{upquote}}{}
\IfFileExists{microtype.sty}{% use microtype if available
  \usepackage[]{microtype}
  \UseMicrotypeSet[protrusion]{basicmath} % disable protrusion for tt fonts
}{}
\makeatletter
\@ifundefined{KOMAClassName}{% if non-KOMA class
  \IfFileExists{parskip.sty}{%
    \usepackage{parskip}
  }{% else
    \setlength{\parindent}{0pt}
    \setlength{\parskip}{6pt plus 2pt minus 1pt}}
}{% if KOMA class
  \KOMAoptions{parskip=half}}
\makeatother
\usepackage{xcolor}
\IfFileExists{xurl.sty}{\usepackage{xurl}}{} % add URL line breaks if available
\IfFileExists{bookmark.sty}{\usepackage{bookmark}}{\usepackage{hyperref}}
\hypersetup{
  pdftitle={Final Report, A Guide to the Architecture Document},
  pdfauthor={Course staff},
  colorlinks=true,
  linkcolor={RoyalBlue},
  filecolor={Maroon},
  citecolor={Blue},
  urlcolor={Blue},
  pdfcreator={LaTeX via pandoc}}
\urlstyle{same} % disable monospaced font for URLs
\usepackage[a4paper,margin=2.4cm]{geometry}
\usepackage{longtable,booktabs,array}
\usepackage{calc} % for calculating minipage widths
% Correct order of tables after \paragraph or \subparagraph
\usepackage{etoolbox}
\makeatletter
\patchcmd\longtable{\par}{\if@noskipsec\mbox{}\fi\par}{}{}
\makeatother
% Allow footnotes in longtable head/foot
\IfFileExists{footnotehyper.sty}{\usepackage{footnotehyper}}{\usepackage{footnote}}
\makesavenoteenv{longtable}
\usepackage{graphicx}
\makeatletter
\def\maxwidth{\ifdim\Gin@nat@width>\linewidth\linewidth\else\Gin@nat@width\fi}
\def\maxheight{\ifdim\Gin@nat@height>\textheight\textheight\else\Gin@nat@height\fi}
\makeatother
% Scale images if necessary, so that they will not overflow the page
% margins by default, and it is still possible to overwrite the defaults
% using explicit options in \includegraphics[width, height, ...]{}
\setkeys{Gin}{width=\maxwidth,height=\maxheight,keepaspectratio}
% Set default figure placement to htbp
\makeatletter
\def\fps@figure{htbp}
\makeatother
\setlength{\emergencystretch}{3em} % prevent overfull lines
\providecommand{\tightlist}{%
  \setlength{\itemsep}{0pt}\setlength{\parskip}{0pt}}
\setcounter{secnumdepth}{-\maxdimen} % remove section numbering
\usepackage{newunicodechar}
\usepackage{textcomp}
\usepackage{amssymb}
\newunicodechar{°}{\textdegree}
\newunicodechar{₂}{\textsubscript{2}}
\newunicodechar{→}{\ensuremath{\rightarrow}}
\newunicodechar{≤}{\ensuremath{\leq}}
\newunicodechar{≥}{\ensuremath{\geq}}
\newunicodechar{≈}{\ensuremath{\approx}}
\newunicodechar{✓}{\ensuremath{\checkmark}}
\newunicodechar{–}{--}
\newunicodechar{—}{---}
\usepackage{etoolbox}
\AtBeginEnvironment{longtable}{\footnotesize}
\AtBeginEnvironment{tabular}{\footnotesize}
\usepackage{fancyhdr}
\pagestyle{fancy}\fancyhf{}
\rhead{\small D7065E\,--\,Architecture Document}
\lhead{\small Final Report -- Section Guide}
\cfoot{\thepage}
\renewcommand{\headrulewidth}{0.4pt}
\usepackage{titlesec}
\titleformat{\section}{\Large\bfseries}{\thesection}{0.6em}{}
\titleformat{\subsection}{\large\bfseries}{\thesubsection}{0.5em}{}
\usepackage{microtype}
\setlength{\emergencystretch}{2em}
\ifLuaTeX
  \usepackage{selnolig}  % disable illegal ligatures
\fi

\title{Final Report, A Guide to the Architecture Document}
\usepackage{etoolbox}
\makeatletter
\providecommand{\subtitle}[1]{% add subtitle to \maketitle
  \apptocmd{\@title}{\par {\large #1 \par}}{}{}
}
\makeatother
\subtitle{D7065E, what each section must contain, and the ideas
behind it}
\author{Course staff}
\date{}

\begin{document}
\maketitle

{
\hypersetup{linkcolor=black}
\setcounter{tocdepth}{3}
\tableofcontents
}
\begin{quote}
\textbf{How to use this guide.} This document explains \emph{how to
write the architecture document} that forms the core of the final report: what each section is for, what it must contain, and the ideas
(MBSE, C4, viewpoints) behind them. It is the companion to the worked
report example, which shows a complete, filled-in version at the grade-5
standard. Read this guide to understand \emph{why}; read the example to
see \emph{what good looks like}. The architecture document is a
\textbf{living document}: start it once the Week 2 proposal is approved
and refine it as the project is built, so the document and the code evolve
together. The final version is submitted with the report and should
match what was actually built.
\end{quote}

\hypertarget{what-the-architecture-document-is}{%
\section{1. What the architecture document
is}\label{what-the-architecture-document-is}}

The architecture document is the \textbf{contract between the design and the implementation}, and it is the main written evidence an
examiner reads. It is where the three kinds of thinking the course is built to develop come together:

\begin{itemize}
\tightlist
\item
  \textbf{Architectural thinking}, modelling the system precisely
  before building it, and justifying every design choice against the alternatives that were rejected.
\item
  \textbf{Systems thinking}, reasoning about the whole feedback loop,
  not isolated parts: how data flows, where computation belongs, and how
  a failure in one component ripples through the rest.
\item
  \textbf{Critical thinking}, judging from evidence how well the
  system works, finding where it breaks, weighing trade-offs, and being
  honest about its weaknesses.
\end{itemize}

A small, well-architected, well-justified, critically-evaluated system
scores higher than a large one that ``just works'' but cannot be
explained or defended. The document is graded across four dimensions
(A--D in the grading criteria); most of the architecture grade is
decided by the quality of the reasoning, not by the number of features.

\hypertarget{the-ideas-behind-the-document}{%
\section{2. The ideas behind the
document}\label{the-ideas-behind-the-document}}

Before the section-by-section guidance, it is worth understanding the
ideas the whole document rests on. Getting these straight is what
separates a report that lists technologies from one that demonstrates
understanding.

\hypertarget{model-based-systems-engineering-lightly}{%
\subsection{2.1 Model-Based Systems Engineering,
lightly}\label{model-based-systems-engineering-lightly}}

The course is grounded in \textbf{Model-Based Systems Engineering
(MBSE)}, but it is easy to mistake the notation for the point, so be
clear about what it means here.

\textbf{The substance of MBSE is architectural and systems thinking}, reasoning about requirements, components, interfaces, behaviour, and the
feedback loop \emph{before} code is written, so design errors surface
while they are still cheap to fix. Errors found during design cost
minutes; the same errors found after deployment cost days, or, in a
building-control system, worse. That early reasoning is the goal, and
it is what this document captures.

\textbf{SysML and heavyweight MBSE tooling are just one way to write
that thinking down, and this course does not use them.} Instead the course uses a few precise, lightweight artefacts, C4 diagrams drawn as code (e.g.\
with D2), a requirements table, interface contracts, and a couple of behavioural
diagrams, kept in the repository next to the code. Using lightweight
artefacts is \emph{not} a departure from MBSE: it keeps MBSE's core (the
thinking) and sheds the notation overhead. What is assessed is the quality of the architectural, systems, and critical reasoning and the ability to justify the decisions, not fluency with a modelling
language.

The single most important MBSE habit is \textbf{traceability}: every
requirement should connect to a design element that satisfies it and a
test that verifies it. If a requirement owns no component, the design
has a gap; if it has no test, the requirement cannot be enforced. This
``requirement → design → validation'' thread runs from the requirements table through the architecture sections to the test plan, and it is the
backbone of a strong document.

\hypertarget{the-c4-model}{%
\subsection{2.2 The C4 model}\label{the-c4-model}}

\textbf{C4} is a lightweight way to draw software architecture, created
by Simon Brown specifically because UML felt too heavy for most teams.
The name comes from its \textbf{four levels of zoom}, each adding
detail. The key idea is to draw the system at the level of detail the reader needs, never everything at once.

\begin{longtable}[]{@{}
  >{\raggedright\arraybackslash}p{(\columnwidth - 6\tabcolsep) * \real{0.1346}}
  >{\raggedright\arraybackslash}p{(\columnwidth - 6\tabcolsep) * \real{0.1154}}
  >{\raggedright\arraybackslash}p{(\columnwidth - 6\tabcolsep) * \real{0.2885}}
  >{\raggedright\arraybackslash}p{(\columnwidth - 6\tabcolsep) * \real{0.4615}}@{}}
\toprule
\begin{minipage}[b]{\linewidth}\raggedright
Level
\end{minipage} & \begin{minipage}[b]{\linewidth}\raggedright
Name
\end{minipage} & \begin{minipage}[b]{\linewidth}\raggedright
What it shows
\end{minipage} & \begin{minipage}[b]{\linewidth}\raggedright
Analogy (Google Maps)
\end{minipage} \\
\midrule
\endhead
1 & \textbf{Context} & The system as a single box, surrounded by its
users and external systems & The country on the globe \\
2 & \textbf{Container} & The separately-deployable units inside
(processes, containers, databases) and how they talk & The cities and
roads \\
3 & \textbf{Component} & The internal parts of \emph{one} container &
The streets within a city \\
4 & \textbf{Code} & Class-level detail (rarely drawn by hand) &
Individual buildings \\
\bottomrule
\end{longtable}

For this course, \textbf{Levels 1 and 2 are what must be produced}, and
they carry most of the architecture grade.

\begin{itemize}
\tightlist
\item
  The \textbf{Context diagram (Level 1)} fixes the \emph{system
  boundary}: what is yours to build versus what is outside it (the
  building manager, occupants, BuildSim, optional external compute).
  Nothing internal is shown. Its job is to make the scope unambiguous.
\item
  The \textbf{Container diagram (Level 2)} opens the box and shows every
  deployable unit, each box is roughly one Docker container, with
  the protocol on every arrow. This is the most important diagram in the
  document, because it should map \textbf{1:1 to the \texttt{docker-compose.yml} and the repository folders}. If the diagram shows six containers and the repository jams everything into
  one folder, the design and the code have parted ways.
\item
  A \textbf{Component diagram (Level 3)} is optional, worth drawing only
  when one container is complex enough to deserve its own picture (for
  example an autonomous service with several internal modules).
\end{itemize}

Beyond the four core levels, C4 has \textbf{supplementary} diagrams that may be used: a \textbf{System Landscape} (the big picture, several systems
and people in one view, no single system in focus), a \textbf{Dynamic}
diagram (how elements collaborate at runtime for one scenario, with
numbered interactions, this is a sequence-style view, in C4 terms),
and a \textbf{Deployment} diagram (what runs where). \textbf{Do not draw
Level 4 (Code) by hand}, C4 itself advises against it, since
class-level diagrams drift from the code immediately and are better
generated by an IDE if ever needed.

\begin{center}\includegraphics[width=0.55\textwidth]{figures/landscape.png}\end{center}

\noindent{\footnotesize\emph{C4 System Landscape, several building
systems, each on BuildSim, overseen by the facilities team; the in-scope
system is highlighted.}}

A \textbf{Level 3 Component} diagram is worth drawing for the one
container complex enough to need it (typically the autonomous service):

\begin{center}\includegraphics[width=\textwidth]{figures/component.png}\end{center}

\noindent{\footnotesize\emph{C4 Component diagram (Level 3), the
internals of one container, its components in light blue.}}

C4 is technology-agnostic and focuses on structure and relationships,
which is exactly the lightweight modelling MBSE asks for.

\hypertarget{the-five-architecture-viewpoints}{%
\subsection{2.3 The five architecture
viewpoints}\label{the-five-architecture-viewpoints}}

No single diagram can serve every reader: a developer needs to see
software components, an operator needs to see machines and networks, a
safety reviewer needs to see failure modes. Trying to show all of that
at once produces an unreadable tangle. The solution is \textbf{multiple
viewpoints}, each a lens that filters out everything not relevant to one
concern, and each catching a kind of error the others hide.

A complete architecture description uses five viewpoints, and they map onto the sections of the report:

\begin{longtable}[]{@{}
  >{\raggedright\arraybackslash}p{(\columnwidth - 4\tabcolsep) * \real{0.1774}}
  >{\raggedright\arraybackslash}p{(\columnwidth - 4\tabcolsep) * \real{0.3387}}
  >{\raggedright\arraybackslash}p{(\columnwidth - 4\tabcolsep) * \real{0.4839}}@{}}
\toprule
\begin{minipage}[b]{\linewidth}\raggedright
Viewpoint
\end{minipage} & \begin{minipage}[b]{\linewidth}\raggedright
Question it answers
\end{minipage} & \begin{minipage}[b]{\linewidth}\raggedright
Where it lives in the report
\end{minipage} \\
\midrule
\endhead
\textbf{Context} & What interacts with the system? & C4 context
diagram \\
\textbf{Functional} & What are the parts, and how do they connect? & C4
container diagram + responsibilities \\
\textbf{Information} & What data exists, and how does it flow? & Data
models, interfaces, pipeline \\
\textbf{Behavioural} & What happens, in what order, for a scenario? &
Sequence and state diagrams \\
\textbf{Deployment} & What runs where? & Deployment diagram + component
placement \\
\bottomrule
\end{longtable}

A correct functional view can still hide a broken deployment view (two
components that look connected might need a network link that does not
exist); a clean data model can hide a behavioural timing bug that only a
sequence diagram exposes. Covering all five avoids leaving a
class of error undiscovered until the demo.

\hypertarget{communication-patterns}{%
\subsection{2.4 Communication patterns}\label{communication-patterns}}

Most architectural mistakes at this scale come from choosing the wrong
way for components to talk. Many patterns are not needed, the right one must be picked for each link and justify it. The three that
matter here:

\begin{itemize}
\tightlist
\item
  \textbf{Request--response (REST over HTTP).} One component asks
  another a specific question and waits. Simple, tightly coupled, no
  buffering. Right for \textbf{commands and queries}, setting an
  actuator's state and confirming it, registering a device, reading a
  value. Wrong for high-frequency streams and for one-to-many fan-out.
\item
  \textbf{Publish--subscribe (MQTT, Kafka).} Publishers send to a
  \emph{topic} without knowing who listens; a broker routes to
  subscribers. Loosely coupled, buffers when a consumer is briefly down.
  Right for \textbf{sensor data distribution}, where one sensor feeds
  many consumers (the database, the autonomous service, the dashboard).
\item
  \textbf{WebSocket.} A persistent two-way connection so the server can
  push to the client without polling. Right for \textbf{live
  dashboards}; BuildSim already uses it to stream the 3D view.
\end{itemize}

The rule of thumb: \textbf{REST when one component asks another and
waits; pub/sub when one source feeds many; WebSocket when the server
pushes to a browser.} State which is used on each link, and why.

\hypertarget{deployment-and-component-placement}{%
\subsection{2.5 Deployment and component
placement}\label{deployment-and-component-placement}}

In this course everything runs on the laptop, so deployment is less
about \emph{where in the world} a component runs and more about
\textbf{how the system is composed from independent processes}. The
deployment viewpoint answers: what runs as its own process, how those
processes are isolated and packaged, how they communicate, and how the
system behaves when one fails. The concerns that matter:

\begin{itemize}
\tightlist
\item
  \textbf{Process isolation and containerization.} Each component runs
  as its own containerized process, so it can be started, stopped, and
  restarted independently, and a fault in one does not take down the
  others.
\item
  \textbf{Communication.} Processes talk only through defined interfaces
  (Section 2.4), never shared memory, which is what lets them be
  developed and deployed independently.
\item
  \textbf{Resource allocation.} Components have different needs, storage for the pipeline, CPU for the autonomous service, little for
  the dashboard, and separate processes let each be provisioned
  independently.
\item
  \textbf{Fault isolation.} If the dashboard crashes, the autonomous
  service keeps running. This is the operational payoff of one
  responsibility per process.
\end{itemize}

The \textbf{edge--cloud continuum} is still worth understanding as part
of the syllabus: a real building would run the control loop on a local
\emph{edge} server so it keeps working when the internet is down, and
use the \emph{cloud} or a GPU server only for slow work off the critical
path (such as overnight model training). In this course that whole
continuum collapses onto one laptop, but the deployment section
should still reason about which components would belong at the edge and
which could be pushed to external compute, and why.

\hypertarget{how-to-justify-a-decision}{%
\subsection{2.6 How to justify a
decision}\label{how-to-justify-a-decision}}

Most of the architecture grade comes from \emph{justification}, and
there is a simple pattern that turns an assertion into a decision:

\begin{quote}
\textbf{``X was chosen because \ldots{} (tie it to a requirement); the alternative was Y; the trade-off accepted is Z.''}
\end{quote}

``MQTT was used'' is an assertion. ``MQTT was used rather than direct REST
because three consumers need the same readings and the broker buffers
them if one is briefly down (NFR-02); the trade-off is one more process
to run and monitor'' is a decision. A decision table with no rejected
alternatives and no trade-offs reads as a list of assumptions, write
the alternatives down even when the choice feels obvious.

\hypertarget{section-by-section-guide}{%
\section{3. Section-by-section guide}\label{section-by-section-guide}}

The rest of this guide walks through each section of the architecture
document. Replace the \textbf{bracketed guidance} with real content;
the illustrative tables in the worked example show the level of detail
expected.

\hypertarget{summary}{%
\subsection{3.1 Summary}\label{summary}}

A short overview so a reader knows, in one glance, what the system is
and where to find things: project title, team, use case, repository, and
version. Include a brief \textbf{``changes since the proposal''} note, what is different from the Week 2 proposal and \emph{why}. The examiner explicitly looks for this, because it shows the design evolved
through reasoning rather than by accident.

\hypertarget{use-case-and-building-context-buildsim}{%
\subsection{3.2 Use case and building context
(BuildSim)}\label{use-case-and-building-context-buildsim}}

Describe the \textbf{use case} and how it connects to \textbf{BuildSim}.
BuildSim is the \emph{shared state} of the cyber-physical system: it
owns the building model (rooms, equipment), exposes a REST API to
register devices and read/write their state, and streams the 3D view
over WebSocket. BuildSim is not modified; it is instrumented. State explicitly how the processes use it, the simulator writes the sensor
values into BuildSim and reads actuator state back (it is the only
process that touches the environment), sensor processes read those values
from BuildSim and publish them onward, actuator processes write their
state to BuildSim, and the dashboard rides BuildSim's WebSocket. Name the actors (building manager,
occupants), tell the concrete scenario the demo will show, and explain
why the use case matters. This section is the \emph{context viewpoint}
in prose.

\hypertarget{requirements}{%
\subsection{3.3 Requirements}\label{requirements}}

State what the system must do as \textbf{testable} statements. Each
requirement gets a unique \textbf{ID}, a \textbf{type}, Functional
(what it does), Non-functional (how well: latency, reliability,
recovery), or Regulatory (an external rule), a \textbf{priority}
(Must / Should / Could), a \textbf{measurable acceptance criterion}, and
the \textbf{test} that verifies it. Non-functional requirements matter
most for grading: they are what the design decisions get justified
against. Quantify everything, ``the system should be fast'' cannot be
tested; ``a decision every 60 s'' can. This table is the root of the
traceability thread from Section 2.1.

\hypertarget{architecture}{%
\subsection{3.4 Architecture}\label{architecture}}

\hypertarget{context-diagram-c4-level-1}{%
\subsubsection{Context diagram (C4 Level
1)}\label{context-diagram-c4-level-1}}

Show the whole system as one box, surrounded by the people and external
systems it touches. Nothing internal. Describe each external actor in a
sentence.

\begin{center}\includegraphics[width=0.64\textwidth]{figures/context.png}\end{center}

\noindent{\footnotesize\emph{C4 System Context (Level 1). Person and software-system elements; the system in scope is highlighted, external systems are grey; every arrow is a directed relationship labelled with its purpose and protocol.}}

\hypertarget{container-diagram-c4-level-2}{%
\subsubsection{Container diagram (C4 Level
2)}\label{container-diagram-c4-level-2}}

Open the box: every deployable unit (≈ one Docker container) and how
they communicate, with the protocol on each arrow. This should map 1:1 to the \texttt{docker-compose.yml}.

\begin{center}\includegraphics[height=0.88\textheight]{figures/container.png}\end{center}

\noindent{\footnotesize\emph{C4 Container diagram (Level 2). Each container shows its [technology] and responsibility; the in-scope system is the dashed boundary, external systems (BuildSim) are grey; every arrow is a directed relationship labelled with its protocol.}}

\hypertarget{component-responsibilities-and-design-decisions}{%
\subsubsection{Component responsibilities and design
decisions}\label{component-responsibilities-and-design-decisions}}

\textbf{Every process should have a single primary responsibility.} List
every container with the \textbf{single job} it owns. For each component,
explain not only \emph{what} it does but \textbf{why it exists as a
separate process} rather than being merged into another component, the
``one job per process'' principle is what makes the system testable and
resilient, and it is exactly what the oral
examination probes. Then give the \textbf{design-decisions table}: for
each significant choice, the option chosen, the alternative(s)
considered, the requirement it serves, and the trade-off accepted (the
pattern from Section 2.6). This table carries most of the architecture
grade.

\hypertarget{interfaces-and-communication}{%
\subsection{3.5 Interfaces and
communication}\label{interfaces-and-communication}}

Specify each interface precisely enough that \textbf{another team could
implement the system from this section alone}. State which communication pattern (Section 2.4) is used on each link and why. Give
the \textbf{message schema} for every boundary crossing, field names,
types, and units, and list each endpoint or topic with its payload
and semantics (status codes, QoS, rate). Ambiguity here is the most
expensive kind of bug: two components that look connected but disagree
on the message format will work in isolation and fail when they meet.

\hypertarget{simulating-the-sensor-values}{%
\subsection{3.6 Simulating the sensor
values}\label{simulating-the-sensor-values}}

Because BuildSim provides the rooms but not their physics, \textbf{the simulator that generates realistic sensor values must be built}. Be explicit
about two things:

\begin{itemize}
\tightlist
\item
  \textbf{The physical processes.} How each sensor value is produced, for example a heat-balance update for temperature (envelope loss, body
  heat, HVAC), and a mass-balance for CO₂ that rises with occupancy and
  falls with ventilation. Crucially, \textbf{actuator commands must feed
  back into the model} so the control loop closes: raising the fan
  lowers CO₂, which the next reading reflects.
\item
  \textbf{The movement of people (occupancy).} Occupancy usually drives
  everything else, heat, CO₂, lighting, so model it explicitly: a
  weekday schedule (arrivals 07:00--09:00, a lunch dip, departures
  16:00--18:00), meeting-room bookings, empty weekends, and some
  randomness so no two days are alike.
\end{itemize}

\textbf{Approaches and examples.} Use a physics-based model (the
equations above), data replay (feed a public dataset through the sensors, for example the UCI \emph{Occupancy Detection} set or
\emph{ASHRAE} building-energy data), or a hybrid (real data with an
injected event). The model need not be physically perfect; it must have
enough structure, trends, daily cycles, genuine responses to
actuators, that the autonomous service reacts to meaningful data. Be
honest about its limitations.

\hypertarget{autonomous-services-and-data-engineering}{%
\subsection{3.7 Autonomous services and data
engineering}\label{autonomous-services-and-data-engineering}}

Describe the \textbf{autonomous service}, the component that turns
data into action. Any approach is acceptable (rule-based, optimisation,
ML, LLM, or a mix); choose the simplest that meets the requirement and
justify it. \textbf{A more sophisticated AI technique does not
automatically result in a higher grade.}
Cover the approach and why it fits the latency budget, the inputs it
uses, the action it produces, how it is integrated, how it is evaluated,
and its failure/safety behaviour (what happens when it is wrong, slow,
or unavailable, and the fallback that keeps the building safe).

Then describe the \textbf{data pipeline}: ingestion (is the broker a
buffer?), storage (a time-series DB for the hot path, a data lake for
history, and why), processing and feature engineering, and how
\textbf{data quality} is handled (missing, stale, or duplicate readings). A
good decision on bad data is still a bad outcome, so data quality is
graded as seriously as the decision logic.

\includegraphics[width=0.98\textwidth,height=\textheight]{figures/pipeline.png}

\hypertarget{behaviour}{%
\subsection{3.8 Behaviour}\label{behaviour}}

Static diagrams show structure; these show \textbf{behaviour over time},
which is where timing and failure bugs hide. Provide a \textbf{dynamic
diagram} (C4's name for a sequence-style view of one scenario, with
numbered interactions among the containers) and a \textbf{state machine}
for at least one process whose lifecycle matters (a sensor or the
autonomous service), including what happens when the network drops and
how it recovers. Every transition in the state machine is a piece of code that must be written.

\begin{center}\includegraphics[width=0.92\textwidth]{figures/dynamic.png}\end{center}

\noindent{\footnotesize\emph{C4 Dynamic diagram, numbered runtime
interactions among the containers for one scenario.}}

\begin{center}\includegraphics[width=\textwidth]{figures/state.png}\end{center}

\hypertarget{deployment-and-component-placement-1}{%
\subsection{3.9 Deployment and component
placement}\label{deployment-and-component-placement-1}}

Show \textbf{where each component runs and justify the placement}. Draw
a deployment diagram with each process and where it runs (in this
course, the laptop), then discuss \textbf{latency, fault isolation,
resource requirements, maintainability, and operational considerations}
(Section 2.5). Everything here runs on one machine, so most of this is
about how the processes are composed, isolated, and communicate; if external compute such as a GPU server is used, explain why it is needed and
how it affects the architecture. Explain \textbf{why each component is
deployed separately and what would be lost if components were merged}.
This is the view previous reports most often forget.

\begin{center}\includegraphics[width=\textwidth]{figures/deployment.png}\end{center}

\hypertarget{test-plan}{%
\subsection{3.10 Test plan}\label{test-plan}}

Explain \textbf{how the system is tested and how evidence is collected},
not every individual test case. Every requirement in Section 3.3 must
have a verification approach and a piece of evidence (a log, a
measurement, a screenshot). Include at least unit tests, one end-to-end
scenario test, and, for the higher grades, fault-injection tests
(kill a process, send corrupt data) and quantitative measurements. This
table is the validation end of the traceability thread.

\hypertarget{evaluation-and-experiments}{%
\subsection{3.11 Evaluation and
experiments}\label{evaluation-and-experiments}}

Present the \textbf{evidence} about how well the system works and where
it breaks. Report the metrics measured (latency, recovery time,
false-alarm rate, energy). For the higher grades, include
\textbf{performance plots and stress experiments}, a
response-over-time plot, an energy comparison against a baseline, a
scaling experiment that finds the breaking point, plus an analysis of
the \textbf{trade-offs} observed and an \textbf{honest critique} of the weaknesses. ``It works'' is a grade-3 statement; ``here is where it
fails, why, and what would be changed'' is a grade-5 statement. The
worked example shows three such plots produced with \texttt{pgfplots}.

\hypertarget{dashboard-risks-and-reflection}{%
\subsection{3.12 Dashboard, risks and
reflection}\label{dashboard-risks-and-reflection}}

Describe what the \textbf{dashboard} visualises (sensor readings,
actuator states, decisions, alerts) and how it is built (BuildSim's
session API or a separate dashboard). Close with \textbf{risks,
assumptions, and a critical reflection}: the things that could undermine
the system, the assumptions the design depends on, and an honest account of its limitations and what would be done next. Being explicit
here is a sign of engineering maturity, not weakness.

\hypertarget{checklist-before-you-submit}{%
\section{4. Checklist before submission}\label{checklist-before-you-submit}}

\begin{itemize}
\tightlist
\item[$\square$]
  Every requirement has an ID, an acceptance criterion, and a test
  (traceability holds).
\item[$\square$]
  C4 context \textbf{and} container diagrams are present and consistent with the \texttt{docker-compose.yml}.
\item[$\square$]
  Every significant design choice is justified against a named
  alternative, with a trade-off.
\item[$\square$]
  Interfaces are specified precisely enough to implement from.
\item[$\square$]
  The autonomous service and the data pipeline are both described and
  justified.
\item[$\square$]
  At least one sequence diagram and one state machine are included.
\item[$\square$]
  The deployment diagram shows where each component runs and justifies
  the placement.
\item[$\square$]
  The evaluation reports measured evidence, finds a breaking point, and
  critiques the design honestly.
\item[$\square$]
  Both team members can explain the whole system, the oral
  examination is individual.
\end{itemize}

\end{document}
